% Right Source, Wrong Number — 4-page workshop paper.
% Self-contained preamble on purpose: compiles on Overleaf or arXiv with no
% external .sty. For a Climate Change AI submission, swap \documentclass for
% their template and delete the geometry/titling block below.
\documentclass[10pt,a4paper]{article}
\usepackage[utf8]{inputenc}
\usepackage[T1]{fontenc}
\usepackage[margin=2.2cm]{geometry}
\usepackage{booktabs}
\usepackage{microtype}
\usepackage{url}
\usepackage[colorlinks=true,allcolors=blue]{hyperref}
\usepackage{titlesec}
\titlespacing*{\section}{0pt}{1.2ex}{0.6ex}
\titlespacing*{\subsection}{0pt}{1.0ex}{0.4ex}
\setlength{\parskip}{0.35em}
\setlength{\parindent}{0pt}
\pagestyle{empty}

\title{\vspace{-2.2em}\textbf{Right Source, Wrong Number:\\Measuring Language-Model Recall of Greenhouse-Gas Emission Factors}\vspace{-0.6em}}
\author{Jeremiah Say\\\small GreenCalculus \quad \texttt{jeremiah@greencalculus.com}}
\date{\vspace{-1em}\small September 2026}

\begin{document}
\maketitle
\vspace{-1.5em}

\begin{abstract}
\small
Greenhouse-gas emission factors are an unusually clean probe for factual recall: each has one correct value, traceable to a named cell in a published document. We put 467 such questions to five frontier language models with no tools and no web access. Unaided accuracy ranged from 41.9\% to 65.2\% within 10\% of the published value. The more consequential result is the failure mode. When these models cited a source they named the \emph{correct} publisher 82--94\% of the time, yet in 31.6--58.8\% of answers that correct attribution accompanied an incorrect number. A wrong value with no provenance is caught in review; a wrong value carrying the right publisher's name is not. In a paired arm on 90 questions, giving two models a sourced lookup tool moved accuracy from 37.1\% to 98.7\% and from 50.0\% to 100.0\%, at 1.0--1.6 tool calls per question, while simultaneously converting refusals into correct answers. We release the questions, every raw model response, the scoring harness, a log of six bugs found in our own scorer, a canary string, and a pre-registered held-out split.
\end{abstract}

\vspace{-0.6em}
\section{Motivation}

Emission factors convert activity into emissions: kilowatt-hours into kilograms of CO\textsubscript{2}e, tonne-kilometres of freight into a carbon figure. They are the arithmetic underneath corporate disclosure, and they are increasingly requested from language models. They also make a good factual-recall benchmark for three reasons. Each factor has a single correct value; that value is traceable to an exact cell in a public document, so ``correct'' is not a matter of judgement; and the values are numerous, jurisdiction-specific and periodically revised, which defeats memorisation of a small canon.

The stakes are not academic. These numbers end up in regulatory filings. Our interest is practical rather than adversarial: we wanted to know whether a model can be trusted to supply one, and if not, what specifically goes wrong.

\vspace{-0.3em}
\section{The benchmark}

\textbf{Questions.} 467 questions drawn deterministically from a sourced corpus at data version 2026.187, spanning 45 categories, 75 publishers and 50 countries. Categories include grid electricity, fuel combustion, well-to-tank, freight, refrigerants, construction materials, spend-based input-output factors, waste and agriculture. Questions are phrased as a practitioner asks (``What is the greenhouse-gas emission factor for UK grid electricity, per kWh?'') rather than as database keys, so the task is recall rather than lookup-by-identifier. Only factors under licences permitting republication were used, so the answer key ships with the benchmark.

\textbf{Ground truth} is the published value together with its unit, gas, publisher, source document and the exact cell reference within it.

\textbf{Scoring.} An answer is correct if it falls within 10\% of the published value after \emph{dimensional reconciliation}: an answer in kg\,CO\textsubscript{2}/GJ is compared against a truth in tonne\,C/GJ as the same physical quantity, handling mass, energy, volume and length conversions, percent-versus-fraction, and carbon-to-CO\textsubscript{2} in both directions. Without this, a correct answer of 63\,kg\,CO\textsubscript{2}/GJ scores as a 366{,}179\% error against 0.0172\,tonne\,C/GJ.

Scoring yields three outcomes, not two. Where units cannot be reconciled mechanically --- a local-currency answer against a US-dollar truth, or a value given with no denominator --- the answer is marked \emph{unscoreable} and removed from the denominator rather than counted wrong. This is the softest part of the method and the reason we report the scored sample size beside every percentage.

\textbf{Citation scoring.} We separately record whether an answer named a source and whether that source was the correct one. Citations count only on answers that actually supplied a number: a refusal that name-drops a publisher (``consult the DEFRA tables'') is not an attribution.

\vspace{-0.3em}
\section{Results}

\begin{table}[h]
\centering
\small
\setlength{\tabcolsep}{4.5pt}
\begin{tabular}{lrrrrrr}
\toprule
Model & Answered & Scored & Within 10\% & Cited & Correct publisher & \textbf{Right src, wrong no.} \\
\midrule
Gemini 3.1 Pro    &  77 &  46 & 65.2\% & 12.0\% & 83.9\% & 35.3\% \\
Grok 4.6          & 144 &  66 & 62.1\% & 21.6\% & 82.2\% & 31.6\% \\
GPT-5.5           & 326 & 256 & 58.2\% & 53.7\% & 94.0\% & 41.0\% \\
Claude Opus 5     & 430 & 315 & 45.7\% & 71.7\% & 89.9\% & 54.1\% \\
Gemini 3.6 Flash  & 404 & 315 & 41.9\% & 71.7\% & 91.9\% & 58.8\% \\
\bottomrule
\end{tabular}
\vspace{-0.3em}
\caption{\small Unaided performance on 467 questions. ``Answered'' counts non-refusals; ``Scored'' excludes unscoreable units. ``Correct publisher'' is conditional on having cited one. The final column is the share of scored answers pairing a correct publisher with an out-of-tolerance value.}
\end{table}
\vspace{-1.2em}

\textbf{Accuracy trades against coverage.} The two most accurate models are the two that answer least. Gemini 3.1 Pro declined 226 of 467 questions and Grok 4.6 declined 313; Claude Opus 5 declined 11. Read the accuracy column alone and caution looks like competence.

\textbf{The failure mode.} Models are good at attribution and bad at values. Conditional on citing anything, they named the right publisher 82--94\% of the time --- but in 31.6\% to 58.8\% of scored answers that correct publisher appeared beside an incorrect number. The models are not confabulating a source to dress up a guess; they know where the number lives and misremember the number. Crucially, the tone is identical either way: we found no signal in the response that distinguishes a correct value from an incorrect one.

This matters for review workflows. An unsourced figure invites checking. A figure carrying ``DEFRA 2026, Table 1'' does not.

\textbf{A sourced lookup closes the gap.} On a paired subset of 90 questions, two models were given two read-only lookup tools and left to decide whether to use them.

\begin{table}[h]
\centering
\small
\setlength{\tabcolsep}{6pt}
\begin{tabular}{lrrr}
\toprule
Model & Unaided & With lookup & Mean tool calls \\
\midrule
Claude Opus 5 & 37.1\% ($n$=62) & \textbf{98.7\%} ($n$=79) & 1.62 \\
GPT-5.5       & 50.0\% ($n$=46) & \textbf{100.0\%} ($n$=79) & 1.00 \\
\bottomrule
\end{tabular}
\vspace{-0.3em}
\caption{\small Paired comparison on identical question ids. $n$ is the scored subset.}
\end{table}
\vspace{-1.2em}

Accuracy and coverage moved together, which is unusual: GPT-5.5 answered 62 of 90 unaided and 87 with the tools, so refusals became correct answers rather than accuracy being purchased with silence.

\vspace{-0.3em}
\section{Two explanations we tested and rejected}

\textbf{Corpus size.} The intuition that finely-subdivided sources are harder to memorise is not supported. Correlating a publisher's corpus size against its accuracy, across the 24 publishers with at least nine scored answers, gives $r = -0.28$; the coefficient ranges from $-0.25$ to $-0.34$ depending on the model pool and inclusion threshold, and at $n = 24$ none of these reaches $p < .05$ (which would require $|r| > 0.40$). The scatter is the clearer argument: a 3-row source scores 55.6\%, a 1{,}424-row source 50.7\%, and a 2{,}451-row source 17.4\%. Size does not order them.

\textbf{Vintage confusion.} An older vintage of one publisher scored 11\% against the current year's 51\%, suggesting models quote the wrong edition. Reading the individual answers killed it: the sample is five questions on unusually obscure aggregates, and none of the answers resembles another vintage of the same factor.

What remains is a pattern we can describe but have not isolated. Reliable categories are old, physically constrained and reprinted widely (fuel properties 93.9\%, waste treatment defaults 100\%); unreliable ones are recent, jurisdiction-specific and published once (CBAM defaults 27.3\%, NGFS scenario prices 9.1\%). At publisher level the contrast is sharper still: the European Commission's CBAM default dataset scores 0\% across every question asked of it. That is consistent with training-data frequency, which we cannot measure from outside. We state it as a pattern rather than a mechanism because two plausible explanations have already failed.

\vspace{-0.3em}
\section{Limitations}

Single domain; results may not transfer to other factual-recall tasks. Point-in-time for the model versions named, September 2026. The 10\% tolerance is a choice, and categories near that boundary would move under a different one. Unscoreable answers were dropped rather than counted wrong --- 28 to 110 per model --- which is the softest methodological decision here. Two of five models have small scored samples (46 and 66) because they declined most questions, so their percentages carry real uncertainty.

Our scorer has been wrong six times; all six are documented, and \emph{four of them understated} the result rather than flattering it. An independent hand-audit of 45 answers agreed with the automated figure to within about one point --- one check, not a proof.

\textbf{Conflict of interest.} The author sells access to emission-factor data, including the lookup tools used in the paired arm. That is precisely why the questions, every raw model response, the ground truth and the scoring code are public: the result should be checkable rather than trusted.

\vspace{-0.3em}
\section{Availability and contamination control}

Questions, all raw model responses, the scorer and the bug log are at \url{https://github.com/greencalculus/greencalculus-benchmark}, archived at \href{https://doi.org/10.5281/zenodo.22692277}{10.5281/zenodo.22692277} and on the HuggingFace Hub as \texttt{greencalculus/emission-factor-benchmark}.

Publishing a benchmark with its answer key degrades it. Two mitigations ship with the release. A \textbf{canary string} is embedded in the question file itself, so it travels if that file is scraped alone; a model that reproduces it was trained on this benchmark and its score here is memorisation. And a \textbf{pre-registered held-out split} of 136 further questions --- same generator, same corpus, disjoint by construction --- is withheld, with only its size, spread and SHA-256 published in advance, so a later release can be verified rather than quietly substituted. From the next evaluation onward we report public and held-out scores side by side; a materially better public score will be reported as memorisation, not accuracy.

\end{document}
